%!TEX TS-program = xelatex
\documentclass[11pt, a4paper]{article}

\usepackage[fontset=macnew]{ctex}
% 修复 macOS 下缺少华文仿宋 (STFangsong) 导致代码环境中调用等宽字体时报错的问题
\renewcommand{\CJKttdefault}{zhkai}


\usepackage{amsmath, amssymb, amsthm}
\usepackage{graphicx}
\usepackage{hyperref}
\usepackage{bookmark}
\usepackage{enumitem}
\usepackage[ruled,vlined]{algorithm2e}

\newtheorem{problem}{题} 
\renewcommand{\proofname}{\textbf{解答}}
\usepackage[margin=1in]{geometry}

\title{算法设计与分析 作业二}
\author{xxxxxxxxx xxx}
\date{\today}

\begin{document}

\maketitle

\begin{problem}[2.7]
\end{problem}

\begin{proof}
    \leavevmode
    \begin{enumerate}[label=(\arabic*)]
    \item 首先使用 $O(n \log n)$ 的复杂度将原数组排序，随后只需要顺序统计，如果有比当前元素更多的新元素，就替换，最后统计最多元素的个数是否大于 $\frac{n}{2}$ 。
    \item 首先可以用 $O(n)$ 的时间找出中位数，随后只需要统计中位数的出现次数是否大于 $\frac{n}{2}$ 即可，时间复杂度为 $O(n)$ 。
    \item 考虑摩尔投票，我们选定一个candidate，然后不同的元素两两抵消，如果当前元素等于candidate，就count++，否则count--，当count==0的时候就换candidate，最后再用candidate去验证它的出现次数，时间复杂度为 $O(n)$ 。
    \end{enumerate}
\end{proof}

\begin{problem}[2.12]
\end{problem}

\begin{proof}
    先对数组 $B$ 排序，然后对于 $A$ 中的不同元素，使用二分查找查找其是否在数组 $B$ 中，若是，则放入数组 $C$ ；时间复杂度为 $O(n \log \log n)$ ，伪代码如下：
    
    \begin{algorithm}[H]
    \caption{求数组交集}
    \KwIn{数组 $A[1..n]$ 和 $B[1..m]$}
    \KwOut{数组 $C$ ，使得 $C = A \cap B$}
    
    $C \leftarrow \emptyset$\;
    $L \leftarrow \text{Sort}(B)$ \tcp*{对 $B$ 排序}
    \For{$i \leftarrow 1$ \KwTo $n$}{
        $x \leftarrow A[i]$\;
        $j \leftarrow \text{BinarySearch}(L, x)$ \tcp*{搜索 $x$ ，返回序标 $j$ 否则返回 0}
        \If{$j > 0$}{
            $C \leftarrow C \cup \{x\}$\;
        }
    }
    \end{algorithm}
\end{proof}

\begin{problem}[2.15]
\end{problem}

\begin{proof}
    \leavevmode
    \begin{enumerate}[label=(\arabic*)]
        \item 对于算法 $A$ ，每次遍历需要 $O(n)$ 次，最坏遍历 $O(\sqrt{n})$ 次，因此最坏情况下的时间复杂度为 $O(n \sqrt{n})$ 。

        对于算法 $B$ ，排序的时间复杂度为 $O(n \log n)$ ，总最坏时间复杂度为 $O(n \log n)$ 。
        \item 考虑使用快速选择算法，线性时间内找出第 $i$ 大元素 $x$ 。

        再遍历数组，收集大于等于 $x$ 的元素并排序，时间复杂度为 $O(n+n+i \log i) = O(n)$ 。
        
        \begin{algorithm}[H]
        \caption{找出数组中最大的 $i$ 个元素并排序}
        \KwIn{数组 $A[1..n]$ ，选取个数 $i$}
        \KwOut{最大 $i$ 个元素的降序数组 $S$}
        
        $x \leftarrow \text{QuickSelect}(A, i)$\;
        $S \leftarrow \emptyset$\;
        $count \leftarrow 0$\;
        \For{$k \leftarrow 1$ \KwTo $n$}{
            \If{$A[k] \ge x$}{
                $S \leftarrow S \cup \{A[k]\}$\;
                $count \leftarrow count + 1$\;
                \If{$count == i$}{
                    \textbf{break}\;
                }
            }
        }
        $S \leftarrow \text{Sort}(S)$\;
        \Return{$S$}\;
        \end{algorithm}
    \end{enumerate}
\end{proof}

\begin{problem}[2.18]
\end{problem}

\begin{proof}
    跟2.15题的(2)相似，先计算出每个点到原点的平方距离 $d_i = x_i^2 + y_i^2$ ；然后仍然使用快速选择算法，在 $O(n)$ 时间内找到距离第 $k = \lfloor\sqrt{n}\rfloor$ 小的元素 $k\_dist$ ；
    
    接下来遍历所有点，把距离小于等于 $k\_dist$ 的点收集起来。最后对收集到的 $\lfloor\sqrt{n}\rfloor$ 个点进行降序排序，并按要求从远到近输出它们的标号。
    
    由于只需对 $\lfloor\sqrt{n}\rfloor$ 个点排序，排序的时间复杂度为 $O(\sqrt{n}\log(\sqrt{n})) \approx O(\sqrt{n}\log n)$ 。总时间复杂度为 $O(n) + O(n) + O(\sqrt{n}\log n) = O(n)$ 。\\
    伪代码如下：
    
    \begin{algorithm}[H]
    \caption{找距离原点最近的 $\lfloor\sqrt{n}\rfloor$ 个点}
    \KwIn{$n$ 个点的坐标 $(x_1, y_1), \dots , (x_n, y_n)$}
    \KwOut{距离最近的 $\lfloor\sqrt{n}\rfloor$ 个点的标号，按距离从远到近输出}
    
    $D \leftarrow \text{大小为 } n \text{ 的空数组}$\;
    \For{$i \leftarrow 1$ \KwTo $n$}{
        $D[i] \leftarrow (x_i^2 + y_i^2, i)$ \tcp*{记录距离的平方以及原下标}
    }
    
    $k \leftarrow \lfloor\sqrt{n}\rfloor$\;
    $x \leftarrow \text{QuickSelect}(D, k)$\; \tcp*{基于距离找到第 $k$ 小的元素}
    
    $S \leftarrow \emptyset$\;
    $count \leftarrow 0$\;
    \For{$i \leftarrow 1$ \KwTo $n$}{
        \If{$D[i].\text{dist} \le x.\text{dist}$}{
            $S \leftarrow S \cup \{D[i]\}$\;
            $count \leftarrow count + 1$\;
            \If{$count = k$}{
                \textbf{break}\;
            }
        }
    }
    
    $S \leftarrow \text{SortDesc}(S)$\; \tcp*{对这 $k$ 个点按距离从大到小排序}
    \For{$p \in S$}{
        \text{Print}(p.\text{index})\;
    }
    \end{algorithm}
\end{proof}

\begin{problem}[3.3]
\end{problem}

\begin{proof}
    相当于是一个01背包问题，时间复杂度最坏为 $O(nD)$ ，伪代码如下：

    \begin{algorithm}[H]
    \caption{货柜存储最大收益算法 (基于动态规划)}
    \KwIn{货柜数量 $n$，库房总长度 $D$，长度数组 $l[1\dots n]$，收益数组 $v[1\dots n]$}
    \KwOut{最大存储收益}
    \tcc{初始化长度为 D+1 的数组，全部置为 0}
    初始化数组 $dp[0 \dots D] \gets 0$\; 
    \For{$i \gets 1$ \KwTo $n$}{
        \tcc{从后向前遍历以保证每个货柜只被放入一次}
        \For{$j \gets D$ \textbf{downto} $l[i]$}{
            $dp[j] \gets \max(dp[j], dp[j - l[i]] + v[i])$\;
        }
    }   
    \Return $dp[D]$\;
    \end{algorithm}
\end{proof}

\begin{problem}[3.8]
\end{problem}

\begin{proof}
    可以视为取值类型的01背包，采用动态规划即可，时间复杂度最坏为 $O(nN)$ ，伪代码如下：

    \begin{algorithm}[H]
    \caption{集合等和划分算法 (基于动态规划)}
    \KwIn{正整数集合 $A = \{a_1, a_2, \dots, a_n\}$，元素总和 $N$}
    \KwOut{布尔值，表示是否能划分成和相等的两个子集}
    
    \If{$N \pmod 2 \neq 0$}{
        \Return \textbf{false}\;
    }
    
    $T \leftarrow N / 2$\;
    \tcc{初始化长度为 T+1 的布尔数组，仅 dp[0] 为 true}
    初始化数组 $dp[0 \dots T] \leftarrow \textbf{false}$\;
    $dp[0] \leftarrow \textbf{true}$\;
    
    \For{$i \leftarrow 1$ \KwTo $n$}{
        \For{$j \leftarrow T$ \textbf{downto} $a_i$}{
            $dp[j] \gets dp[j] \lor dp[j - a_i]$\;
        }
    }
    \Return $dp[T]$\;
    \end{algorithm}
\end{proof}

\begin{problem}[3.10]
\end{problem}

\begin{proof}
    这就是二维的01背包，只需要在转移的时候多跑一个体积的维度即可，时间复杂度最坏为 $O(nWV)$ ，伪代码如下：
    
    \begin{algorithm}[H]
    \caption{二维费用 0-1 背包问题}
    \KwIn{物品数量 $n$，最大重量 $W$，最大体积 $V$，物品的重量 $w[1 \dots n]$、体积 $c[1 \dots n]$ 和价值 $v[1 \dots n]$}
    \KwOut{最大价值}
    
    \tcc{初始化所有状态为0}
    初始化二维数组 $dp[0 \dots W][0 \dots V] \leftarrow 0$\;
    
    \For{$i \leftarrow 1$ \KwTo $n$}{
        \tcc{重量和体积均需要倒序遍历以保证物品只被装入一次}
        \For{$j \leftarrow W$ \textbf{downto} $w[i]$}{
            \For{$k \leftarrow V$ \textbf{downto} $c[i]$}{
                $dp[j][k] \gets \max(dp[j][k], dp[j - w[i]][k - c[i]] + v[i])$\;
            }
        }
    }
    \Return $dp[W][V]$\;
    \end{algorithm}
\end{proof}

\begin{problem}[3.17]
\end{problem}

\begin{proof}
    典型的数字三角形问题，只需要算出每个点的dp值，然后从最后一层贪心一一回溯即可，伪代码如下：
    
    \begin{algorithm}[H]
    \caption{数字三角形最小路径和算法 (自底向上)}
    \KwIn{数字三角形的行数 $n$，各节点的值 $A[1\dots n][1\dots n]$}
    \KwOut{最小路径和，以及对应的路径}
    
    \tcc{初始化动态规划数组 $dp$ 和路径记录数组 $choice$}
    $dp \leftarrow \text{大小为 } n \times n \text{ 的数组}$\;
    $choice \leftarrow \text{大小为 } n \times n \text{ 的数组}$\;
    
    \tcc{初始化最后一行}
    \For{$j \leftarrow 1$ \KwTo $n$}{
        $dp[n][j] \leftarrow A[n][j]$\;
    }
    
    \tcc{从倒数第二行开始向上递推}
    \For{$i \leftarrow n-1$ \textbf{downto} $1$}{
        \For{$j \leftarrow 1$ \KwTo $i$}{
            \If{$dp[i+1][j] < dp[i+1][j+1]$}{
                $dp[i][j] \leftarrow A[i][j] + dp[i+1][j]$\;
                $choice[i][j] \leftarrow 0$ \tcp*{0 表示向下走左边}
            }
            \Else{
                $dp[i][j] \leftarrow A[i][j] + dp[i+1][j+1]$\;
                $choice[i][j] \leftarrow 1$ \tcp*{1 表示向下走右边}
            }
        }
    }
    
    \tcc{根据 choice 数组还原路径}
    $curr\_j \leftarrow 1$\;
    $path \leftarrow \text{空列表}$\;
    \For{$i \leftarrow 1$ \KwTo $n-1$}{
        $\text{把 } A[i][curr\_j] \text{ 加入 } path$\;
        $curr\_j \leftarrow curr\_j + choice[i][curr\_j]$\;
    }
    $\text{把 } A[n][curr\_j] \text{ 加入 } path$\;
    
    \Return $dp[1][1], path$\;
    \end{algorithm}
\end{proof}

\vspace{1em}

\end{document}